\documentclass{article}
\usepackage[vietnamese]{babel}
\usepackage{amsmath,amssymb}
\usepackage{geometry}
\geometry{left=2cm,right=2cm,top=2cm,bottom=2cm}
\usepackage{pgfplots}
\pgfplotsset{compat=1.18}
\documentclass[12pt,a4paper]{article}
\usepackage[vietnamese]{babel}
\usepackage{amsmath,amssymb}
\usepackage{geometry}
\geometry{left=2cm,right=2cm,top=2cm,bottom=2cm}
\usepackage{graphicx}
\usepackage{tikz}
\usetikzlibrary{calc}

\begin{document}

\usetikzlibrary{calc}

\begin{titlepage}
    % Tạo khung viền đôi cho trang bìa
    \begin{tikzpicture}[overlay, remember picture]
        \draw[line width=1.5pt, blue!50!black] 
            ($ (current page.north west) + (1.2cm,-1.2cm) $) 
            rectangle 
            ($ (current page.south east) + (-1.2cm,1.2cm) $);
        \draw[line width=0.5pt, blue!50!black] 
            ($ (current page.north west) + (1.35cm,-1.35cm) $) 
            rectangle 
            ($ (current page.south east) + (-1.35cm,1.35cm) $);
    \end{tikzpicture}

    \centering
    \vspace*{3cm}
    

    {\Large \textbf{DỰ ÁN}\\[0.5cm]}
    {\Large \textbf{ỨNG DỤNG HÌNH HỌC CỦA TÍCH PHÂN}}

    \vfill

    % Phần thông tin sinh viên thực hiện ở giữa trang (hoặc phía dưới một chút)
    \textbf{Sinh viên thực hiện}\\
    \vspace{0.3}
    Trượng Ngụy Minh Triết\\
    MSV: 23S1010046

    \vfill
    {\normalsize Huế, 6/2026}
\end{titlepage}

\section*{1. KHÁI NIỆM VÀ TÍCH CHẤT CỦA TÍCH PHÂN}

\subsection*{a. Khái niệm tích phân}

Cho $f(x)$ là hàm số liên tục trên đoạn $[a, b]$. Nếu $F(x)$ là một nguyên hàm của $f(x)$ thì hiệu số $F(b) - F(a)$ được gọi là tích phân từ $a$ đến $b$ của hàm số $f(x)$, kí hiệu là:

\[
\int_{a}^{b} f(x) \, dx
\]

\paragraph{Chú ý.}

\begin{enumerate}
    \item[1)] Hiệu $F(b) - F(a)$ thường được kí hiệu là $\left. F(x) \right|_a^b$.
    
    Như vậy
    \[
    \int_a^b f(x)\mathrm{d}x = \left. F(x) \right|_a^b.
    \]

    \item[2)] Ta gọi $\displaystyle\int_a^b$ là dấu tích phân, $a$ là cận dưới, $b$ là cận trên, $f(x)\mathrm{d}x$ là \textit{biểu thức dưới dấu tích phân} và $f(x)$ là hàm số dưới dấu tích phân.

    \item[3)] Trong trường hợp $a = b$ hoặc $a > b$, ta quy ước:
    \[
    \int_a^a f(x)\mathrm{d}x = 0; \quad \int_a^b f(x)\mathrm{d}x = -\int_b^a f(x)\mathrm{d}x.
    \]
\end{enumerate}

\noindent\textbf{Ví dụ.} Tính:
\[
\int_{-1}^{3} x^2 \, \mathrm{d}x
\]

\noindent\textbf{Giải:}
\[
\int_{-1}^{3} x^2 \, \mathrm{d}x = \left. \frac{x^3}{3} \right|_{-1}^{3} = \frac{1}{3} \left[ 3^3 - (-1)^3 \right] = \frac{28}{3}.
\]

\subsection*{b. Tính chất của tích phân}

\begin{enumerate}
    \item $\int_{a}^{a} f(x) \, dx = 0$
    \item $\int_{a}^{b} f(x) \, dx = -\int_{b}^{a} f(x) \, dx$
    \item $\int_{a}^{b} k \cdot f(x) \, dx = k \cdot \int_{a}^{b} f(x) \, dx \quad (k \text{ là hằng số})$
    \item $\int_{a}^{b} [f(x) \pm g(x)] \, dx = \int_{a}^{b} f(x) \, dx \pm \int_{a}^{b} g(x) \, dx$
    \item $\int_{a}^{c} f(x) \, dx = \int_{a}^{b} f(x) \, dx + \int_{b}^{c} f(x) \, dx \quad (\forall b \in [a, c])$
    \item Nếu $f(x) \ge 0, \forall x \in [a, b]$ thì $\int_{a}^{b} f(x) \, dx \ge 0$.
\end{enumerate}

\subsection*{c. Các công thức nguyên hàm, tích phân}

\paragraph{1. Bảng công thức nguyên hàm cơ bản:}
\[

\begin{center}
\renewcommand{\arraystretch}{2.2}
\begin{tabular}{|p{7.5cm}|p{7.5cm}|}
    \hline
    \centering \textbf{Nguyên hàm của các hàm số sơ cấp thường gặp} & 
    \centering \textbf{Nguyên hàm của các hàm số hợp \newline (dưới đây $u = u(x)$)} \tabularnewline \hline
    
    $\displaystyle\int \mathrm{d}x = x + C$ & 
    $\displaystyle\int \mathrm{d}u = u + C$ \tabularnewline \hline
    
    $\displaystyle\int x^{\alpha}\mathrm{d}x = \frac{x^{\alpha+1}}{\alpha+1} + C \quad (\alpha \neq -1)$ & 
    $\displaystyle\int u^{\alpha}\mathrm{d}u = \frac{u^{\alpha+1}}{\alpha+1} + C \quad (\alpha \neq -1)$ \tabularnewline \hline
    
    $\displaystyle\int \frac{\mathrm{d}x}{x} = \ln|x| + C \quad (x \neq 0)$ & 
    $\displaystyle\int \frac{\mathrm{d}u}{u} = \ln|u| + C \quad (u = u(x) \neq 0)$ \tabularnewline \hline
    
    $\displaystyle\int e^x \mathrm{d}x = e^x + C$ & 
    $\displaystyle\int e^u \mathrm{d}u = e^u + C$ \tabularnewline \hline
    
    $\displaystyle\int a^x \mathrm{d}x = \frac{a^x}{\ln a} + C \quad (0 < a \neq 1)$ & 
    $\displaystyle\int a^u \mathrm{d}u = \frac{a^u}{\ln a} + C \quad (0 < a \neq 1)$ \tabularnewline \hline
    
    $\displaystyle\int \cos x \, \mathrm{d}x = \sin x + C$ & 
    $\displaystyle\int \cos u \, \mathrm{d}u = \sin u + C$ \tabularnewline \hline
    
    $\displaystyle\int \sin x \, \mathrm{d}x = -\cos x + C$ & 
    $\displaystyle\int \sin u \, \mathrm{d}u = -\cos u + C$ \tabularnewline \hline
    
    $\displaystyle\int \frac{\mathrm{d}x}{\cos^2 x} = \int (1+\tan^2 x)\mathrm{d}x = \tan x + C$ & 
    $\displaystyle\int \frac{\mathrm{d}u}{\cos^2 u} = \int (1+\tan^2 u)\mathrm{d}u = \tan u + C$ \tabularnewline \hline
    
    $\displaystyle\int \frac{\mathrm{d}x}{\sin^2 x} = \int (1+\cot^2 x)\mathrm{d}x = -\cot x + C$ & 
    $\displaystyle\int \frac{\mathrm{d}u}{\sin^2 u} = \int (1+\cot^2 u)\mathrm{d}u = -\cot u + C$ \tabularnewline \hline
    
    $\displaystyle\int \frac{\mathrm{d}x}{2\sqrt{x}} = \sqrt{x} + C \quad (x > 0)$ & 
    $\displaystyle\int \frac{\mathrm{d}u}{2\sqrt{u}} = \sqrt{u} + C \quad (u > 0)$ \tabularnewline \hline
\end{tabular}
\end{center}

\paragraph{2. Hàm số lượng giác:}
\[

\begin{center}
\renewcommand{\arraystretch}{2.2} % Khoảng cách dòng thoáng
\begin{tabular}{|p{5cm}|p{5cm}|p{6.5cm}|}
    \hline
    \centering \textbf{Nguyên hàm của hàm số sơ cấp} & 
    \centering \textbf{Nguyên hàm của hàm số hợp ($u=u(x)$)} & 
    \centering \textbf{Nguyên hàm của hàm số hợp ($u=ax+b;\ a\neq 0$)} \tabularnewline \hline
    
    $\displaystyle\int \sin x \, \mathrm{d}x = -\cos x + C$ & 
    $\displaystyle\int \sin u \, \mathrm{d}u = -\cos u + C$ & 
    $\displaystyle\int \sin(ax+b)\,\mathrm{d}x = -\frac{1}{a}\cos(ax+b) + C$ \tabularnewline \hline
    
    $\displaystyle\int \cos x \, \mathrm{d}x = \sin x + C$ & 
    $\displaystyle\int \cos u \, \mathrm{d}u = \sin u + C$ & 
    $\displaystyle\int \cos(ax+b)\,\mathrm{d}x = \frac{1}{a}\sin(ax+b) + C$ \tabularnewline \hline
    
    $\displaystyle\int \tan x \, \mathrm{d}x = -\ln|\cos x| + C$ & 
    $\displaystyle\int \tan u \, \mathrm{d}u = -\ln|\cos u| + C$ & 
    $\displaystyle\int \tan(ax+b)\,\mathrm{d}x = -\frac{1}{a}\ln|\cos(ax+b)| + C$ \tabularnewline \hline
    
    $\displaystyle\int \cot x \, \mathrm{d}x = \ln|\sin x| + C$ & 
    $\displaystyle\int \cot u \, \mathrm{d}u = \ln|\sin u| + C$ & 
    $\displaystyle\int \cot(ax+b)\,\mathrm{d}x = \frac{1}{a}\ln|\sin(ax+b)| + C$ \tabularnewline \hline
    
    $\displaystyle\int \frac{1}{\sin^2 x}\,\mathrm{d}x = -\cot x + C$ & 
    $\displaystyle\int \frac{1}{\sin^2 u}\,\mathrm{d}u = -\cot u + C$ & 
    $\displaystyle\int \frac{1}{\sin^2(ax+b)}\,\mathrm{d}x = -\frac{1}{a}\cot(ax+b) + C$ \tabularnewline \hline
    
    $\displaystyle\int \frac{1}{\cos^2 x}\,\mathrm{d}x = \tan x + C$ & 
    $\displaystyle\int \frac{1}{\cos^2 u}\,\mathrm{d}u = \tan u + C$ & 
    $\displaystyle\int \frac{1}{\cos^2(ax+b)}\,\mathrm{d}x = \frac{1}{a}\tan(ax+b) + C$ \tabularnewline \hline
    
    $\displaystyle\int \frac{1}{\sin x}\,\mathrm{d}x = \ln\left|\tan\frac{x}{2}\right| + C$ & 
    $\displaystyle\int \frac{1}{\sin u}\,\mathrm{d}u = \ln\left|\tan\frac{u}{2}\right| + C$ & 
    $\displaystyle\int \frac{\mathrm{d}x}{\sin(ax+b)} = \frac{1}{a}\ln\left|\tan\frac{ax+b}{2}\right| + C$ \tabularnewline \hline
    
    $\displaystyle\int \frac{1}{\cos x}\,\mathrm{d}x = \ln\left|\tan\left(\frac{x}{2}+\frac{\pi}{4}\right)\right|+C$ & 
    $\displaystyle\int \frac{1}{\cos u}\,\mathrm{d}u = \ln\left|\tan\left(\frac{u}{2}+\frac{\pi}{4}\right)\right|+C$ & 
    $\displaystyle\int \frac{\mathrm{d}x}{\cos(ax+b)} = \frac{1}{a}\ln\left|\tan\left(\frac{ax+b}{2}+\frac{\pi}{4}\right)\right|+C$ \tabularnewline \hline
\end{tabular}
\end{center}
\paragraph{3. Bảng công thức nguyên hàm nâng cao:}
\[
\begin{align*}
    &\int \sqrt{a^2 - x^2} \, \mathrm{d}x = \frac{x\sqrt{a^2 - x^2}}{2} + \frac{a^2}{2}\arcsin\frac{x}{a} + C \\
    &\int \frac{\mathrm{d}x}{\sin(ax + b)} = \frac{1}{a}\ln\left|\tan\frac{ax + b}{2}\right| + C \\
    &\int e^{ax}\cos bx \, \mathrm{d}x = \frac{e^{ax}(a\cos bx + b\sin bx)}{a^2 + b^2} + C \\
    &\int e^{ax}\sin bx \, \mathrm{d}x = \frac{e^{ax}(a\sin bx - b\cos bx)}{a^2 + b^2} + C \\
    &\int \frac{\mathrm{d}x}{a^2 + x^2} = \frac{1}{a}\arctan\frac{x}{a} + C \\
    &\int \frac{\mathrm{d}x}{a^2 - x^2} = \frac{1}{2a}\ln\left|\frac{a + x}{a - x}\right| + C \\
    &\int \frac{\mathrm{d}x}{\sqrt{x^2 + a^2}} = \ln\left(x + \sqrt{x^2 + a^2}\right) + C \\
    &\int \frac{\mathrm{d}x}{\sqrt{a^2 - x^2}} = \arcsin\frac{x}{|a|} + C \\
    &\int \frac{\mathrm{d}x}{x\sqrt{x^2 - a^2}} = \frac{1}{a}\arccos\left|\frac{a}{x}\right| + C \\
    &\int \frac{\mathrm{d}x}{x\sqrt{x^2 + a^2}} = -\frac{1}{a}\ln\left|\frac{a + \sqrt{x^2 + a^2}}{x}\right| + C \\
    &\int \ln(ax + b) \, \mathrm{d}x = \left(x + \frac{b}{a}\right)\ln(ax + b) - x + C
\end{align*}

\subsection*{d. Các phương pháp tính tích phân}

\begin{itemize}
    \item \textbf{Tính tích phân bằng định nghĩa, tính chất và bảng nguyên hàm cơ bản}: Sử dụng trực tiếp bảng nguyên hàm các hàm số cơ bản, các tính chất tuyến tính của tích phân và định nghĩa tích phân để tính toán.
    
    \item \textbf{Phương pháp tích phân từng phần}: 
    Sử dụng công thức $\int_{a}^{b} u(x)v'(x)\,\mathrm{d}x = u(x)v(x)\Big|_a^b - \int_{a}^{b} v(x)u'(x)\,\mathrm{d}x$ hoặc viết gọn là $\int_{a}^{b} u\,\mathrm{d}v = uv\Big|_a^b - \int_{a}^{b} v\,\mathrm{d}u$.
    Các bước thực hiện gồm: viết biểu thức dưới dấu tích phân dưới dạng $u\,\mathrm{d}v$, tính $\mathrm{d}u$ và $v$, sau đó áp dụng công thức. 
    Khi chọn $u$ và $\mathrm{d}v$, đối với dạng đa thức nhân hàm mũ hoặc lượng giác thì đặt $u$ là đa thức; đối với dạng đa thức nhân lô-garit thì đặt $u$ là hàm lô-garit; đối với dạng tích phân chứa $e^{ax}\cos bx$ hoặc $e^{ax}\sin bx$ cần thực hiện tích phân từng phần hai lần liên tiếp để quay lại tích phân ban đầu.

    \item \textbf{Phương pháp đổi biến số}:
    \begin{itemize}
        \item \textit{Phương pháp đổi biến dạng I}: Đặt ẩn phụ $x = u(t)$ sao cho hàm hợp và vi phân được đơn giản hóa, đồng thời đổi cận tương ứng. Trường hợp chứa các căn thức $\sqrt{a^2-x^2}$, $\sqrt{a^2+x^2}$, $\sqrt{x^2-a^2}$ mà không biến đổi được bằng cách khác, ta đặt ẩn phụ bằng các hàm lượng giác thích hợp tương ứng để khử căn.
        \item \textit{Phương pháp đổi biến dạng II}: Đặt ẩn phụ $u = u(x)$, tính vi phân $\mathrm{d}u = u'(x)\mathrm{d}x$ để đưa biểu thức về biến $u$ hoàn toàn và đổi cận từ $a, b$ sang $u(a), u(b)$.
    \end{itemize}

    \item \textbf{Tích phân hàm số phân thức hữu tỉ $\int \frac{P(x)}{Q(x)}\mathrm{d}x$}: 
    Nếu bậc của $P(x) \ge$ bậc của $Q(x)$ thì thực hiện phép chia đa thức. 
    Nếu bậc của $P(x) <$ bậc của $Q(x)$, phân tích thành các phân thức đơn giản tùy theo nghiệm của mẫu số $Q(x)$ như nghiệm đơn, nhân tử bậc hai vô nghiệm, nghiệm bội, hoặc sử dụng phương pháp đồng nhất hệ số đối với tam thức bậc hai ở mẫu.

    \item \textbf{Tích phân các hàm lượng giác}: 
    Sử dụng các công thức biến đổi cơ bản (tích thành tổng, hạ bậc) hoặc các phép đặt ẩn phụ chuyên dụng như đặt $t = \tan\frac{x}{2}$ cho dạng tổng quát $\int \frac{\mathrm{d}x}{a\sin x + b\cos x + c}$; chia tử và mẫu cho $\cos^2 x$ hoặc đặt $t = \tan x$ đối với dạng chứa $\sin^2 x, \cos^2 x, \sin x\cos x$ ở mẫu; hoặc dùng hệ số bất định cho dạng phân thức chứa bậc nhất của sin và cos. 
    Ngoài ra có thể áp dụng tính chẵn lẻ của hàm số đối với $\sin x$ và $\cos x$ để chọn phép đặt $t = \tan x$ ($\cot x$), $t = \cos x$, hoặc $t = \sin x$.

    \item \textbf{Tích phân hàm vô tỉ}: 
    Khử căn thức bằng cách nhân lượng liên hợp, biến đổi thành bình phương đúng, hoặc đặt toàn bộ căn thức (hoặc phần phức tạp) làm ẩn phụ $t$ để làm mất hẳn căn.

    \item \textbf{Tích phân chứa dấu giá trị tuyệt đối}: 
    Lập bảng xét dấu của biểu thức trong dấu giá trị tuyệt đối trên miền lấy tích phân, sau đó chia miền tích phân thành các khoảng nhỏ để phá dấu giá trị tuyệt đối và tính từng phần.

    \item \textbf{Tích phân một số hàm đặc biệt}:
    \begin{itemize}
        \item Hàm số liên tục và lẻ trên đoạn $[-a; a]$ có tích phân xác định bằng $0$: $\int_{-a}^{a} f(x)\,\mathrm{d}x = 0$.
        \item Hàm số liên tục và chẵn trên đoạn $[-a; a]$ có thể rút gọn bằng hai lần tích phân trên nửa đoạn dương: $\int_{-a}^{a} f(x)\,\mathrm{d}x = 2\int_{0}^{a} f(x)\,\mathrm{d}x$.
        \item Dạng tích phân có mẫu chứa hàm mũ $\int_{-\alpha}^{\alpha} \frac{f(x)}{a^x+1}\,\mathrm{d}x = \frac{1}{2}\int_{-\alpha}^{\alpha} f(x)\,\mathrm{d}x$.
        \item Tính chất đối xứng của hàm lượng giác trên $[0; \frac{\pi}{2}]$: $\int_{0}^{\frac{\pi}{2}} f(\sin x)\,\mathrm{d}x = \int_{0}^{\frac{\pi}{2}} f(\cos x)\,\mathrm{d}x$.
    \end{itemize}
\end{itemize}

\section*{2. ỨNG DỤNG TÍCH PHÂN ĐỂ TÍNH DIỆN TÍCH HÌNH PHẲNG}

\subsection*{a. Hình phẳng giới hạn bởi một đồ thị hàm số, trục hoành và hai đường thẳng $x = a, x = b$}

\paragraph{Khái niệm.} 
Diện tích $S$ của hình phẳng giới hạn bởi đồ thị của hàm số $f(x)$ liên tục, trục hoành và hai đường thẳng $x = a, x = b$ ($a < b$), được tính bằng công thức:
\[
S = \int_{a}^{b} |f(x)| \, \mathrm{d}x
\]

\paragraph{Ví dụ 1.} Tính diện tích hình phẳng giới hạn bởi đồ thị của hàm số $y = x^3$, trục hoành và hai đường thẳng $x = 0, x = 2$.

\noindent\textbf{Giải:}  
Diện tích hình phẳng cần tính là:
\[
S = \int_{0}^{2} |x^3| \, \mathrm{d}x = \int_{0}^{2} x^3 \, \mathrm{d}x = \left. \frac{x^4}{4} \right|_{0}^{2} = \frac{2^4}{4} - 0 = 4 - 0 = 4.
\]

\begin{center}
\begin{tikzpicture}
    \begin{axis}[
        xmin=-0.5, xmax=2.5,
        ymin=-1, ymax=9,
        axis lines=middle,
        xlabel=$x$, ylabel=$y$,
        xtick={0,2},
        ytick={2,4,6,8},
        samples=100,
        width=7cm, height=6cm,
    ]
        \addplot[name path=f, blue, domain=0:2] {x^3};
        \addplot[name path=axis, draw=none] {0};
        \addplot[pink!50] fill between[of=f and axis];
        \draw[dashed] (axis cs:2,0) -- (axis cs:2,8);
    \end{axis}
\end{tikzpicture}
\end{center}

\paragraph{Ví dụ 2.} Tính diện tích hình phẳng giới hạn bởi đồ thị hàm số $y = \sin x$, trục hoành và hai đường thẳng $x = 0, x = 2\pi$.

\noindent\textbf{Giải:}  
Diện tích hình phẳng cần tính là:
\[
S = \int_{0}^{2\pi} |\sin x| \, \mathrm{d}x = \int_{0}^{\pi} |\sin x| \, \mathrm{d}x + \int_{\pi}^{2\pi} |\sin x| \, \mathrm{d}x
\]
\[
= \int_{0}^{\pi} \sin x \, \mathrm{d}x + \int_{\pi}^{2\pi} (-\sin x) \, \mathrm{d}x
\]
\[
= \left. -\cos x \right|_{0}^{\pi} + \left. \cos x \right|_{\pi}^{2\pi} = (-\cos \pi + \cos 0) + (\cos 2\pi - \cos \pi) = (1 + 1) + (1 - (-1)) = 4.
\]

\begin{center}
\begin{tikzpicture}
    \begin{axis}[
        xmin=-0.5, xmax=7,
        ymin=-1.3, ymax=1.3,
        axis lines=middle,
        xlabel=$x$, ylabel=$y$,
        xtick={0, 3.1416, 6.2831},
        xticklabels={$0$, $\pi$, $2\pi$},
        ytick={-1,1},
        samples=200,
        width=10cm, height=5cm,
    ]
        \addplot[name path=f, blue, domain=0:6.2831] {sin(deg(x))};
        \addplot[name path=axis, draw=none, domain=0:6.2831] {0};
        \addplot[pink!50] fill between[of=f and axis, domain=0:3.1416];
        \addplot[pink!50] fill between[of=f and axis, domain=3.1416:6.2831];
        \draw[dashed] (axis cs:3.1416,0) -- (axis cs:3.1416,1);
        \draw[dashed] (axis cs:6.2831,0) -- (axis cs:6.2831,-1);
    \end{axis}
\end{tikzpicture}
\end{center}

\paragraph{Luyện tập 1.} Tính diện tích hình phẳng giới hạn bởi parabol $y = x^2 - 4$, trục hoành và hai đường thẳng $x = 0, x = 3$.

\begin{center}
\begin{tikzpicture}
    \begin{axis}[
        xmin=-1.5, xmax=4.5,
        ymin=-5, ymax=6,
        axis lines=middle,
        xlabel=$x$, ylabel=$y$,
        xtick={-1,1,2,3,4},
        ytick={-4,-3,-2,-1,1,2,3,4,5},
        samples=100,
        width=7.5cm, height=7cm,
    ]
        \addplot[name path=f, magenta, domain=-1:3.3] {x^2 - 4};
        \addplot[name path=axis, draw=none, domain=-1:3.3] {0};
        \addplot[orange!60] fill between[of=f and axis, domain=0:3];
        \draw[dashed, thick] (axis cs:3,0) -- (axis cs:3,5);
        \node[magenta, font=\small, rotate=50] at (axis cs:1.2, -1.8) {$y = x^2 - 4$};
    \end{axis}
\end{tikzpicture}
\end{center}

\vspace{0.5cm}

\subsection*{b. Hình phẳng giới hạn bởi hai đồ thị hàm số và hai đường thẳng $x = a, x = b$}

\paragraph{Khái niệm.} 
Diện tích $S$ của hình phẳng giới hạn bởi đồ thị của hai hàm số $f(x), g(x)$ liên tục trên đoạn $[a; b]$ và hai đường thẳng $x = a, x = b$, được tính bằng công thức:
\[
S = \int_{a}^{b} |f(x) - g(x)| \, \mathrm{d}x
\]

\paragraph{Chú ý.} Nếu hiệu $f(x) - g(x)$ không đổi dấu trên đoạn $[a; b]$ thì:
\[
\int_{a}^{b} |f(x) - g(x)| \, \mathrm{d}x = \left| \int_{a}^{b} (f(x) - g(x)) \, \mathrm{d}x \right|
\]

\paragraph{Ví dụ 3.} Tính diện tích hình phẳng giới hạn bởi hai parabol $y = 4 - x^2$, $y = x^2$ và hai đường thẳng $x = -1, x = 1$.

\noindent\textbf{Giải:}  
Diện tích hình phẳng cần tính là:
\[
S = \int_{-1}^{1} |(4 - x^2) - x^2| \, \mathrm{d}x = \int_{-1}^{1} |4 - 2x^2| \, \mathrm{d}x
\]
\[
= \int_{-1}^{1} (4 - 2x^2) \, \mathrm{d}x = \left. \left( 4x - \frac{2}{3}x^3 \right) \right|_{-1}^{1} = \left( 4 - \frac{2}{3} \right) - \left( -4 + \frac{2}{3} \right) = \frac{10}{3} - \left( -\frac{10}{3} \right) = \frac{20}{3}.
\]

\begin{center}
\begin{tikzpicture}
    \begin{axis}[
        xmin=-2.5, xmax=2.5,
        ymin=-1.5, ymax=5.5,
        axis lines=middle,
        xlabel=$x$, ylabel=$y$,
        xtick={-1,1},
        ytick={1,2,3,4},
        samples=100,
        width=8cm, height=7.5cm,
    ]
        \addplot[name path=f, magenta, domain=-2.3:2.3] {4 - x^2};
        \addplot[name path=g, blue, domain=-1.9:1.9] {x^2};
        \addplot[orange!60] fill between[of=f and g, domain=-1:1];
        \draw[dashed, thick] (axis cs:-1,0) -- (axis cs:-1,3);
        \draw[dashed, thick] (axis cs:1,0) -- (axis cs:1,3);
        \node[magenta, font=\small] at (axis cs:-1.5, 3.8) {$y = 4 - x^2$};
        \node[blue, font=\small] at (axis cs:1.3, 3.3) {$y = x^2$};
    \end{axis}
\end{tikzpicture}
\end{center}

\paragraph{Ví dụ 4.} Tính diện tích hình phẳng giới hạn bởi đồ thị hai hàm số $y = \sin x$, $y = \cos x$ và hai đường thẳng $x = 0, x = \frac{\pi}{4}$.

\noindent\textbf{Giải:}  
Diện tích hình phẳng cần tính là:
\[
S = \int_{0}^{\frac{\pi}{4}} |\sin x - \cos x| \, \mathrm{d}x = \int_{0}^{\frac{\pi}{4}} (\cos x - \sin x) \, \mathrm{d}x
\]
\[
= \left. (\sin x + \cos x) \right|_{0}^{\frac{\pi}{4}} = \left( \frac{\sqrt{2}}{2} + \frac{\sqrt{2}}{2} \right) - (0 + 1) = \sqrt{2} - 1.
\]

\begin{center}
\begin{tikzpicture}
    \begin{axis}[
        xmin=-0.2, xmax=1.6,
        ymin=-0.2, ymax=1.3,
        axis lines=middle,
        xlabel=$x$, ylabel=$y$,
        xtick={0.7854},
        xticklabels={$\frac{\pi}{4}$},
        ytick={0.707},
        yticklabels={$\frac{\sqrt{2}}{2}$},
        samples=200,
        width=9cm, height=6cm,
    ]
        \addplot[name path=f, blue, domain=0:1.3] {sin(deg(x))};
        \addplot[name path=g, red, domain=0:1.3] {cos(deg(x))};
        \addplot[orange!60] fill between[of=g and f, domain=0:0.7854];
        \draw[dashed, thick] (axis cs:0.7854,0) -- (axis cs:0.7854,0.707);
        \node[blue, font=\small] at (axis cs:1.2, 1.05) {$y = \sin x$};
        \node[red, font=\small] at (axis cs:1.2, 0.22) {$y = \cos x$};
    \end{axis}
\end{tikzpicture}
\end{center}

\paragraph{Luyện tập 2.} Tính diện tích hình phẳng giới hạn bởi đồ thị của các hàm số $y = \sqrt{x}$, $y = x - 2$ và hai đường thẳng $x = 1, x = 4$.

\section*{3. ỨNG DỤNG TÍCH PHÂN ĐỂ TÍNH THỂ TÍCH VẬT THỂ}

\subsection*{a. Tính thể tích của vật thể}

\paragraph{Công thức tính thể tích của vật thể.} 
Cho một vật thể trong không gian $Oxyz$. Gọi $\mathcal{B}$ là phần vật thể giới hạn bởi hai mặt phẳng vuông góc với trục $Ox$ tại các điểm có hoành độ $x = a, x = b$. Một mặt phẳng vuông góc với trục $Ox$ tại điểm có hoành độ là $x$ cắt vật thể theo mặt cắt có diện tích là $S(x)$. Giả sử $S(x)$ là hàm số liên tục trên đoạn $[a; b]$. Khi đó thể tích $V$ của phần vật thể $\mathcal{B}$ được tính bởi công thức:
\[
V = \int_{a}^{b} S(x) \, \mathrm{d}x
\]

\paragraph{Ví dụ.} Tính thể tích của khối chóp đều có đáy là hình vuông cạnh $L$ và chiều cao là $h$.

\noindent\textbf{Giải:} 
Chọn trục $Ox$ sao cho gốc $O$ trùng với đỉnh của khối chóp và trục đi qua tâm của đáy. Khi đó, đáy của khối chóp nằm trên mặt phẳng vuông góc với $Ox$ tại $x = h$. 
Mỗi mặt phẳng vuông góc với trục $Ox$ tại điểm có hoành độ bằng $x$ ($0 \le x \le h$), cắt khối chóp theo mặt cắt là hình vuông có cạnh là $a$. 
Theo định lí Thales, ta có $\frac{x}{h} = \frac{a}{L}$, suy ra $a = \frac{L}{h}x$. 
Do đó, diện tích của mặt cắt này là $S(x) = a^2 = \frac{L^2}{h^2}x^2$. 

Vậy thể tích của khối chóp này là:
\[
V = \int_{0}^{h} S(x) \, \mathrm{d}x = \int_{0}^{h} \frac{L^2}{h^2}x^2 \, \mathrm{d}x = \left. \frac{L^2}{h^2} \frac{x^3}{3} \right|_{0}^{h} = \frac{1}{3}L^2h.
\]

\begin{center}
\begin{tikzpicture}[scale=0.9, >=stealth]
    % Trục tọa độ
    \draw[->] (-0.5,0) -- (5,0) node[right] {$x$};
    \draw[->] (0,-1.8) -- (0,2.2) node[above] {$y$};
    \fill (0,0) circle (1.5pt) node[below left] {$O$};
    
    \def\h{4}
    \def\px{2.2}
    
    % Tam giác cắt (màu xanh lá)
    \draw[green!60!black, thick] (0,0) -- (\h, 1.5) node[above left] {$P$};
    \draw[green!60!black, thick] (0,0) -- (\h, -1.5);
    \draw[green!60!black, thick] (\h, 1.5) -- (\h, -1.5);
    
    % Đoạn thẳng mặt cắt a tại x (màu đỏ)
    \draw[red, thick] (\px, 0.825) -- (\px, -0.825);
    % Đặt chữ a thấp xuống sát trục Ox
    \node[red, above right, inner sep=1pt] at (\px, 0) {$a$};
    \fill[red] (\px, 0) circle (1.5pt);
    
    % Điểm L trên trục Ox và đặt chữ L sát trục Ox
    \fill[red] (\h, 0) circle (1.5pt);
    \node[red, above right, inner sep=1pt] at (\h, 0) {$L$};
    
    % Các đường gióng kích thước x và h ở phía dưới
    \draw[gray, thin, <->] (0, -1.0) -- (\px, -1.0) node[midway, below] {$x$};
    \draw[gray, thin, <->] (0, -1.5) -- (\h, -1.5) node[midway, below] {$h$};
\end{tikzpicture}
\end{center}

\paragraph{Chú ý.} Bằng ứng dụng của tích phân, người ta chứng minh được thể tích của khối chóp bất kì bằng $\frac{1}{3}$ diện tích mặt đáy nhân với chiều cao của nó.

\vspace{0.5cm}

\subsection*{b. Tính thể tích khối tròn xoay}

\paragraph{Công thức tính thể tích khối tròn xoay.} 
Cho hàm số $f(x)$ liên tục, không âm trên đoạn $[a; b]$. 
Khi quay hình phẳng giới hạn bởi đồ thị hàm số $y = f(x)$, trục hoành và hai đường thẳng $x = a, x = b$ quanh trục hoành, ta được hình khối gọi là \textbf{khối tròn xoay}. 
Khi cắt khối tròn xoay đó bởi một mặt phẳng vuông góc với trục $Ox$ tại điểm $x \in [a; b]$ được một hình tròn có bán kính $f(x)$. Thể tích của khối tròn xoay này là:
\[
V = \pi \int_{a}^{b} f^2(x) \, \mathrm{d}x
\]

\paragraph{Ví dụ.} Tính thể tích khối tròn xoay sinh ra khi quay quanh trục $Ox$ hình phẳng giới hạn bởi đồ thị hàm số $y = \sqrt{x}$, trục hoành và hai đường thẳng $x = 0, x = 1$.

\noindent\textbf{Giải:} 
Thể tích khối tròn xoay cần tính là:
\[
V = \pi \int_{0}^{1} f^2(x) \, \mathrm{d}x = \pi \int_{0}^{1} (\sqrt{x})^2 \, \mathrm{d}x = \pi \int_{0}^{1} x \, \mathrm{d}x = \left. \pi \frac{x^2}{2} \right|_{0}^{1} = \frac{\pi}{2}.
\]

\begin{center}
\begin{tikzpicture}
    \begin{axis}[
        xmin=-0.5, xmax=1.5,
        ymin=-1.2, ymax=1.5,
        axis lines=middle,
        xlabel=$x$, ylabel=$y$,
        xtick={1},
        ytick={},
        samples=100,
        width=8cm, height=6cm,
    ]
        \addplot[name path=f, blue, domain=0:1.2] {sqrt(x)};
        \addplot[name path=axis, draw=none, domain=0:1.2] {0};
        \addplot[orange!60] fill between[of=f and axis, domain=0:1];
        \node[blue, font=\small] at (axis cs:0.7, 1.0) {$y = \sqrt{x}$};
        \draw[dashed, thick] (axis cs:1,0) -- (axis cs:1,1);
    \end{axis}
\end{tikzpicture}
\end{center}

\section*{4. ỨNG DỤNG TÍCH PHÂN ĐỂ TÍNH DIỆN TÍCH MẶT TRÒN XOAY}

Cho hình thang cong giới hạn bởi 
$\begin{cases} 
a \le x \le b \\ 
y = 0 \\ 
y = f(x) 
\end{cases}$ 
với $f \in C^1[a, b]$. Quay hình thang cong này quanh trục $Ox$ thì ta được một vật thể tròn xoay. Khi đó diện tích xung quanh của vật thể được tính theo công thức:

\begin{equation}
S = 2\pi \int_{a}^{b} |f(x)| \sqrt{1 + f'^2(x)} \, \mathrm{d}x
\end{equation}

Tương tự nếu quay hình thang cong 
$\begin{cases} 
c \le y \le d \\ 
x = 0 \\ 
x = \varphi(y) 
\end{cases}$ 
với $\varphi \in C^1[c, d]$, quanh trục $Oy$ thì:

\begin{equation}
S = 2\pi \int_{c}^{d} |\varphi(y)| \sqrt{1 + \varphi'^2(y)} \, \mathrm{d}y
\end{equation}

\bigskip

\noindent\textbf{Ví dụ.} Tính diện tích mặt tròn xoay tạo nên khi quay các đường sau:

\noindent\textbf{a/} $y = \tan x, \, 0 < x \le \frac{\pi}{4}$ quanh trục $Ox$.

\noindent\textbf{b/} $\frac{x^2}{a^2} + \frac{y^2}{b^2} = 1$ quanh trục $Oy \, (a > b)$.

\noindent\textbf{c/} $9y^2 = x(3 - x)^2, \, 0 \le x \le 3$ quanh trục $Ox$.

\bigskip

\noindent\textbf{Lời giải:}

\noindent\textbf{a/} Áp dụng công thức (1) ta có:
\[
S = 2\pi \int_{0}^{\frac{\pi}{4}} \tan x \sqrt{1 + (1 + \tan^2 x)} \, \mathrm{d}x
\]
\[
= 2\pi \int_{0}^{1} t \sqrt{1 + (1 + t^2)^2} \cdot \frac{\mathrm{d}t}{1 + t^2} \quad (\text{đặt } t = \tan x)
\]
\[
= \pi \int_{0}^{1} \frac{\sqrt{1 + (1 + t^2)^2}}{1 + t^2} \, \mathrm{d}(t^2 + 1)
\]
\[
= \pi \int_{1}^{2} \frac{\sqrt{1 + s^2}}{s} \, \mathrm{d}s \quad (\text{đặt } s = 1 + t^2)
\]
\[
= \pi \int_{\sqrt{2}}^{\sqrt{5}} \left(1 + \frac{1}{u^2 - 1}\right) \mathrm{d}u
\]
\[
= \pi \left\{ \sqrt{5} - \sqrt{2} + \frac{1}{2} \left[ \ln \frac{\sqrt{5} - 1}{\sqrt{2} - 1} - \ln \frac{\sqrt{5} + 1}{\sqrt{2} + 1} \right] \right\}
\]

\bigskip

\noindent\textbf{b/} Nhận xét tính đối xứng của miền và áp dụng công thức (2) ta có:
\[
S = 2 \cdot 2\pi \int_{0}^{b} \frac{a}{b} \sqrt{b^2 - y^2} \cdot \sqrt{1 + \left( \frac{a}{b} \frac{y}{\sqrt{b^2 - y^2}} \right)^2} \, \mathrm{d}y
\]
\[
= 4\pi \frac{a}{b} \int_{0}^{b} \sqrt{b^4 + (a^2 - b^2)y^2} \, \mathrm{d}y
\]
\[
= 4\pi \frac{a}{b} \sqrt{a^2 - b^2} \int_{0}^{b} \sqrt{y^2 + \frac{b^4}{a^2 - b^2}} \, \mathrm{d}y
\]
\[
= 4\pi \frac{a}{b} \sqrt{a^2 - b^2} \int_{0}^{b} \sqrt{y^2 + \beta} \, \mathrm{d}y \quad \left( \text{đặt } \beta = \frac{b^4}{a^2 - b^2} \right)
\]
\[
= 4\pi \frac{a}{b} \sqrt{a^2 - b^2} \cdot \frac{1}{2} \left[ y \sqrt{y^2 + \beta} + \beta \ln |y + \sqrt{y^2 + \beta}| \right] \Bigg|_{0}^{b}
\]
\[
\left( \int \sqrt{y^2 + \beta} \, \mathrm{d}y = \frac{1}{2} \left[ y \sqrt{y^2 + \beta} + \beta \ln |y + \sqrt{y^2 + \beta}| \right] \right)
\]

\bigskip

\noindent\textbf{c/} Trước hết:
\[
9y^2 = x(3 - x)^2 \implies 18yy' = 3(3 - x)(1 - x) \implies y' = \frac{(3 - x)(1 - x)}{6y} \implies y'^2 = \frac{(1 - x)^2}{4x}
\]
Nên áp dụng công thức (1) ta có:
\[
S = 2\pi \int_{0}^{3} \frac{\sqrt{x}(3 - x)}{3} \cdot \sqrt{1 + \frac{(1 - x)^2}{4x}} \, \mathrm{d}x = 2\pi \cdot \frac{1}{6} \int_{0}^{3} (3 - x)(1 + x) \, \mathrm{d}x = 3\pi
\]

\section*{5. ỨNG DỤNG TÍCH PHÂN TÍNH ĐỘ DÀI ĐƯỜNG CONG PHẲNG}

\noindent\textbf{a) Trường hợp đường cong $AB$ cho bởi phương trình $y = f(x)$:}
\[
AB \begin{cases} 
y = f(x) \\ 
a \le x \le b \\ 
f \in C^1[a, b] 
\end{cases} 
\quad \text{thì} \quad 
s = \int_{a}^{b} \sqrt{1 + [f'(x)]^2} \, \mathrm{d}x \tag{3}
\]

\bigskip

\noindent\textbf{b) Trường hợp đường cong $AB$ cho bởi phương trình tham số:}
\[
AB \begin{cases} 
x = x(t) \\ 
y = y(t) \\ 
\alpha \le t \le \beta \\ 
x(t), y(t) \in C^1[a, b] \\ 
x'^2(t) + y'^2(t) > 0 \quad \forall t \in [\alpha, \beta] 
\end{cases} 
\quad \text{thì} \quad 
s = \int_{\alpha}^{\beta} \sqrt{[x'(t)]^2 + [y'(t)]^2} \, \mathrm{d}t \tag{4}
\]

\bigskip

\noindent\textbf{c) Trường hợp đường cong $AB$ cho bởi phương trình trong toạ độ cực:}
\[
AB \begin{cases} 
r = r(\varphi) \\ 
\alpha \le \varphi \le \beta \\ 
r(\varphi) \in C^1[\alpha, \beta] 
\end{cases} 
\quad \text{thì} \quad 
s = \int_{\alpha}^{\beta} \sqrt{r^2(\varphi) + r'^2(\varphi)} \, \mathrm{d}\varphi \tag{5}
\]

\bigskip

\noindent\textbf{Ví dụ.} Tính độ dài đường cong:

\noindent\textbf{a/} $y = \ln \frac{e^x + 1}{e^x - 1}$ khi $x$ biến thiên từ $1$ đến $2$.

\noindent\textbf{b/} 
$\begin{cases} 
x = a\left(\cos t + \ln \tan \frac{t}{2}\right) \\ 
y = a \sin t 
\end{cases}$ 
khi $t$ biến thiên từ $\frac{\pi}{3}$ đến $\frac{\pi}{2}$.

\bigskip

\noindent\textbf{Lời giải:}

\noindent\textbf{a/} Ta có:
\[
1 + y'^2(x) = 1 + \left( \frac{e^x}{e^x + 1} - \frac{e^x}{e^x - 1} \right)^2 = \left( \frac{e^{2x} + 1}{e^{2x} - 1} \right)^2
\]
Nên áp dụng công thức (3) ta được:
\[
s = \int_{1}^{2} \frac{e^{2x} + 1}{e^{2x} - 1} \, \mathrm{d}x \;\; \stackrel{(t = e^{2x})}{=} \int_{e^2}^{e^4} \frac{t + 1}{2t(t - 1)} \, \mathrm{d}t = \ln \frac{e^2 + 1}{e^2}
\]

\bigskip

\noindent\textbf{b/} Áp dụng công thức (4) ta có:
\[
x'^2(t) + y'^2(t) = a^2 \cdot \frac{\cos^2 t}{\sin^2 t} \implies s = a \int_{\frac{\pi}{3}}^{\frac{\pi}{2}} \sqrt{\frac{\cos^2 t}{\sin^2 t}} \, \mathrm{d}t = a \ln \sqrt{\frac{2}{\sqrt{3}}}
\]

\section*{6. MỘT SỐ BÀI TẬP ỨNG DỤNG HÌNH HỌC CỦA TÍCH PHÂN}

\noindent\textbf{Bài 1.} Các nhà kinh tế sử dụng đường cong Lorenz để minh hoạ sự phân phối thu nhập trong một quốc gia. Gọi $x$ là đại diện cho phần trăm số gia đình trong một quốc gia và $y$ là phần trăm tổng thu nhập. Đường $y = x$ sẽ đại diện cho một quốc gia mà các gia đình có thu nhập như nhau. Đường cong Lorenz $y = f(x)$ biểu thị phân phối thu nhập thực tế. Diện tích hình phẳng giới hạn bởi hai đường này, với $0 \le x \le 100$, biểu thị ``sự bất bình đẳng về thu nhập'' của một quốc gia. Năm $2005$, đường cong Lorenz của Hoa Kỳ có thể được mô hình hoá bởi hàm số
\[
y = \left(0{,}00061x^2 + 0{,}0218x + 1{,}723\right)^2, \quad 0 \le x \le 100,
\]
trong đó $x$ được tính từ các gia đình nghèo nhất đến giàu có nhất (Theo \textit{R. Larson, Brief Calculus: An Applied Approach, 8th edition, Cengage Learning, 2009}).

\noindent\textbf{Lời giải:}

\noindent\textbf{Đáp án:} $297945768{,}18$.

\noindent Sự bất bình đẳng thu nhập của Hoa Kỳ vào năm 2005 là:
\begin{align*}
    &\int_{0}^{100} \left| \left(0{,}00061x^2 + 0{,}0218x + 1723\right)^2 - x \right| \mathrm{d}x \\
    &= \int_{0}^{100} \left| 0{,}00061^2x^4 + 4{,}7524 \cdot 10^{-4} \cdot x^2 + 1723^2 \right. \\
    &\quad \left. + 2{,}6596 \cdot 10^{-5}x^3 + 2{,}10206x^2 + 74{,}1228x \right| \mathrm{d}x \\
    &= \int_{0}^{100} \left( 0{,}00061^2x^4 + 4{,}7524 \cdot 10^{-4} \cdot x^2 + 1723^2 \right. \\
    &\quad \left. + 2{,}6596 \cdot 10^{-5}x^3 + 2{,}10206x^2 + 74{,}1228x \right) \mathrm{d}x \\
    &= \left( 7{,}442 \cdot 10^{-8}x^5 + \frac{11881}{75} \cdot 10^{-6}x^3 + 1723^2x \right. \\
    &\quad \left. + 6{,}649 \cdot 10^{-6}x^4 + \frac{105103}{150000}x^3 + 37{,}0614x^2 \right)\Bigg|_{0}^{100} \\
    &= 7{,}442 \cdot 10^{-8} \cdot 100^5 + \frac{11881}{75} \cdot 10^{-6} \cdot 100^3 + 1723^2 \cdot 100 \\
    &\quad + 6{,}649 \cdot 10^{-6} \cdot 100^4 + \frac{105103}{150000} \cdot 100^3 + 37{,}0614 \cdot 100^2 \\
    &= 297945768{,}18.
\end{align*}

\noindent\textbf{Bài 2.} Khối chỏm cầu có bán kính $R$ và chiều cao $h$ ($0 < h \le R$) sinh ra khi quay hình phẳng giới hạn bởi cung tròn có phương trình $y = \sqrt{R^2 - x^2}$, trục hoành và hai đường thẳng $x = R - h, x = R$ xung quanh trục $Ox$. Tính thể tích của khối chỏm cầu này.

\begin{center}
\begin{tikzpicture}[scale=1.1, >=stealth]
    % ==================== HÌNH 1 (BÊN TRÁI) ====================
    \begin{scope}[shift={(0,0)}]
        % Trục tọa độ
        \draw[->, thick] (0,0) -- (3.2,0) node[below, cyan] {$x$};
        \draw[->, thick] (0,-2.2) -- (0,2.5) node[left, cyan] {$y$};
        \draw[->, thick] (0,0) -- (-2.2,-1.5) node[above, cyan] {$z$};
        \node[above left, cyan] at (0,0) {$O$};

        % Miền tô vàng chỏm cầu
        \fill[yellow!40] (1.2,0) -- (1.2,{sqrt(2.2^2 - 1.2^2)}) arc (57.02:0:2.2) -- cycle;

        % Nửa đường tròn bán kính R
        \draw[thick, red] (0,2.2) arc (90:-90:2.2);

        % Mặt cắt elip tại x = R - h: nửa ngoài liền, nửa trong đứt
        \draw[thick] (1.2,{sqrt(2.2^2 - 1.2^2)}) arc (90:-90:0.283 and {sqrt(2.2^2 - 1.2^2)});
        \draw[thick, dashed] (1.2,{sqrt(2.2^2 - 1.2^2)}) arc (90:270:0.283 and {sqrt(2.2^2 - 1.2^2)});

        % Các đoạn thẳng chú thích R-h và h trên trục Ox
        \draw[thick, green!60!black] (0,0) -- (1.2,0);
        \draw[thick, orange] (1.2,0) -- (2.2,0);
        \node[above, orange] at (1.7,0) {$h$};
        \node[below, cyan] at (2.2,0) {$R$};
        
        % Chú thích R - h bằng mũi tên 
        \node[green!60!black] (rh) at (0.3,-1.2) {$R-h$};
        \draw[->, gray, thick] (rh) |- (0.6,-0.05);

        % Điểm đỏ trên trục
        \fill[red] (1.2,0) circle (1.5pt);

        % Nhãn hàm số y = sqrt(R^2 - x^2)
        \node[red, right] at (1.8, 1.9) {$y = \sqrt{R^2 - x^2}$};
    \end{scope}

    % ==================== HÌNH 2 (BÊN PHẢI) ====================
    \begin{scope}[shift={(6,0)}]
        % Khối cầu bên ngoài
        \draw[thick, red] (0,0) circle (2);

        % Đường xích đạo
        \draw[thick] (-2,0) arc (180:360:2 and 0.5);
        \draw[thick, dashed] (2,0) arc (0:180:2 and 0.5);

        % Mặt cắt chỏm cầu màu xanh
        \fill[cyan!60] (0,1.2) ellipse ({sqrt(2^2 - 1.2^2)} and 0.35);
        \fill[cyan!60] ({-sqrt(2^2 - 1.2^2)},1.2) arc (143.13:36.87:2) -- cycle;
        \draw[thick] ({-sqrt(2^2 - 1.2^2)},1.2) arc (180:360:{sqrt(2^2 - 1.2^2)} and 0.35);
        \draw[thick, dashed] ({sqrt(2^2 - 1.2^2)},1.2) arc (0:180:{sqrt(2^2 - 1.2^2)} and 0.35);

        % Các điểm tâm
        \fill[red] (0,0) circle (1.8pt);
        \fill[red] (0,1.2) circle (1.8pt);

        % Bán kính R
        \draw[cyan, dashed, thick] (0,0) -- (2,0);
        \node[cyan, above] at (1,0) {$R$};

        % Đo chiều cao h
        \draw[dashed, gray] (0,2) -- (2.1,2);
        \draw[dashed, gray] ({sqrt(2^2 - 1.2^2)},1.2) -- (2.1,1.2);
        \draw[<->, >=stealth, thick, gray] (1.9,1.2) -- (1.9,2);
        \node[orange, right] at (1.9,1.6) {$h$};
    \end{scope}
\end{tikzpicture}
\end{center}

\noindent\textbf{Lời giải:}

\noindent Thể tích khối chỏm cầu là:
\begin{align*}
    V &= \pi \int_{R-h}^{R} \left(R^2 - x^2\right) \mathrm{d}x \\
    &= \pi \left( R^2x - \frac{x^3}{3} \right)\Bigg|_{R-h}^{R} \\
    &= \pi \left[ R^3 - \frac{R^3}{3} - R^2(R-h) + \frac{(R-h)^3}{3} \right] \\
    &= \pi h^2 \left( R - \frac{h}{3} \right).
\end{align*}

\noindent\textbf{Bài 3.} Cho tam giác vuông $OAB$ có cạnh $OA = a$ nằm trên trục $Ox$ và $\widehat{AOB} = \alpha$ $\left(0 < \alpha \le \dfrac{\pi}{4}\right)$. Gọi $\mathcal{B}$ là khối tròn xoay sinh ra khi quay miền tam giác $OAB$ xung quanh trục $Ox$.
\begin{enumerate}
    \item[a)] Tính thể tích $V$ của $\mathcal{B}$ theo $a$ và $\alpha$.
    \item[b)] Tìm $\alpha$ sao cho thể tích $V$ lớn nhất.
\end{enumerate}

\begin{center}
\begin{tikzpicture}[scale=1.2, >=stealth]
    % Hệ trục tọa độ 3 chiều Oxyz
    \draw[->, thick] (-1, 0) -- (5, 0) node[below] {$x$};
    \draw[->, thick] (0, -2.2) -- (0, 2.2) node[left] {$y$};
    \draw[->, thick] (0, 0) -- (-1.5, -1.5) node[below left] {$z$};

    % Các tọa độ điểm
    \coordinate (O) at (0, 0);
    \coordinate (A) at (3.5, 0);
    \coordinate (B) at (3.5, 1.4);

    % Chỉ tô màu vàng đúng phần tam giác OAB ở nửa trên
    \fill[yellow!35] (O) -- (B) -- (A) -- cycle;

    % Vẽ các đường sinh và đường cao
    \draw[dashed, thick] (O) -- (A); 
    \draw[thick] (O) -- (B);       
    \draw[thick] (O) -- (3.5, -1.4); 
    \draw[thick] (A) -- (B);       

    % Đáy hình nón (vẽ hình elip)
    \draw[thick] (3.5, 0) ellipse (0.25cm and 1.4cm);

    % Góc alpha (thu ngắn bán kính lại để không bị nhô ra ngoài)
    \draw (0.9, 0) arc (0:21.8:0.9cm);
    \node at (1.1, 0.22) {$\alpha$};

    % Tên các điểm O, A, B viết rõ ràng
    \node[fill=white, inner sep=1pt] at (O) [above left] {$O$};
    \node[fill=white, inner sep=1pt] at (A) [below right, yshift=-1pt] {$A$};
    \node[fill=white, inner sep=1pt] at (B) [above] {$B$};
\end{tikzpicture}
\end{center}

\noindent\textbf{Lời giải:}

\noindent a) Ta có: $AB = a \tan\alpha$. Khi quay tam giác $AOB$ quanh trục $Ox$ ta được khối nón tròn xoay có bán kính đáy $r = AB = a \tan\alpha$ và chiều cao $h = OA = a$.

\noindent Thể tích khối nón là:
\[
V = \frac{1}{3}\pi r^2 h = \frac{1}{3}\pi \cdot a \cdot a^2\tan^2\alpha = \frac{1}{3}\pi a^3\tan^2\alpha.
\]

\noindent b) Theo a ta có: $V = \frac{1}{3}\pi a^3 \tan^2\alpha$.

\noindent Ta có: $V' = \frac{2}{3}\pi a^3 \tan\alpha \cdot \frac{1}{\cos^2\alpha}$. Với $0 < \alpha \le \dfrac{\pi}{4} \Rightarrow 0 < \tan\alpha < 1$. Do đó, $V' > 0$ nên hàm số $V$ đồng biến trên $\left(0; \dfrac{\pi}{4}\right)$.

\noindent Do đó, $\max\limits_{\left(0;\frac{\pi}{4}\right]} V = V\left(\frac{\pi}{4}\right) = \frac{1}{3}\pi a^3 \tan^2\frac{\pi}{4} = \frac{1}{3}\pi a^3$.

\noindent Vậy giá trị lớn nhất của $V$ là $\dfrac{1}{3}\pi a^3$ khi $\alpha = \dfrac{\pi}{4}$.


\noindent\textbf{Bài 4.} Kiến trúc sư thiết kế một khu sinh hoạt cộng đồng có dạng hình chữ nhật với chiều rộng và chiều dài lần lượt là $60\,\text{m}$ và $80\,\text{m}$. Trong đó, phần được tô màu đậm là sân chơi, phần còn lại để trồng hoa. Mỗi phần trồng hoa có đường biên cong là một phần của parabol với đỉnh thuộc một trục đối xứng của hình chữ nhật và khoảng cách từ đỉnh đó đến trung điểm cạnh tương ứng của hình chữ nhật bằng $20\,\text{m}$ (xem hình minh họa). Diện tích của phần sân chơi là bao nhiêu mét vuông?

\begin{center}
\begin{tikzpicture}
    \begin{axis}[
        xmin=-10, xmax=70,
        ymin=-5, ymax=85,
        axis lines=none,
        samples=100,
        width=6.5cm, height=8cm, % Đã thu nhỏ kích thước lại cho gọn
    ]
        % --- TÔ MÀU PHẦN THÂN CHÍNH Ở GIỮA, ĐỂ TRẮNG 2 ĐẦU PARABOL ---
        \addplot[fill=brown!15, draw=none, domain=0:60] {60 + 0.02222*(x-30)^2} \closedcycle;
        \addplot[fill=white, draw=none, domain=0:60] {20 - 0.02222*(x-30)^2} \closedcycle;

        % --- VẼ VIỀN HÌNH CHỮ NHẬT VÀ ĐƯỜNG CONG PARABOL ---
        \addplot[thick, black] coordinates {(0,0) (60,0) (60,80) (0,80) (0,0)};

        % Đường cong parabol trên và dưới màu cam
        \addplot[thick, orange!80!black, domain=0:60] {60 + 0.02222*(x-30)^2};
        \addplot[thick, orange!80!black, domain=0:60] {20 - 0.02222*(x-30)^2};

        % --- CÁC MŨI TÊN ĐO KHOẢNG CÁCH 20 M KHÍT ---
        \draw[<->, >=stealth, thick] (axis cs:30, 80) -- (axis cs:30, 60);
        \node[fill=white, inner sep=1pt] at (axis cs:30, 70) {\small $20\text{ m}$};

        \draw[<->, >=stealth, thick] (axis cs:30, 20) -- (axis cs:30, 0);
        \node[fill=white, inner sep=1pt] at (axis cs:30, 10) {\small $20\text{ m}$};

    \end{axis}
\end{tikzpicture}
\end{center}

\noindent\textbf{Lời giải:}

\noindent\textbf{Đáp án:} $3200$.

\noindent Đặt hệ trục tọa độ $Oxy$ như hình vẽ.

\begin{center}
\begin{tikzpicture}[scale=0.9, >=stealth]
    % Sử dụng pgfplots để vẽ hệ trục chuẩn xác
    \begin{axis}[
        xmin=-10, xmax=75,
        ymin=-10, ymax=95,
        axis lines=middle,
        xlabel=$x$, ylabel=$y$,
        xtick={30,60},
        xticklabels={$30$, $60$},
        ytick={20,80},
        yticklabels={$20$, $80$},
        samples=100,
        width=8cm, height=9cm,
    ]
        % Hình chữ nhật tổng thể (chiều rộng 60, chiều cao 80)
        \addplot[thick, black] coordinates {(0,0) (60,0) (60,80) (0,80) (0,0)};
        
        % Tô màu phần sân chơi ở giữa (giới hạn bởi 2 parabol)
        \addplot[fill=brown!15, draw=none, domain=0:60] {80 - (-1/45*x^2 + 4/3*x)} \closedcycle;
        \addplot[fill=white, domain=0:60] {-1/45*x^2 + 4/3*x} \closedcycle;

        % Parabola phía dưới và phía trên
        \addplot[thick, orange!80!black, domain=0:60] {-1/45*x^2 + 4/3*x};
        \addplot[thick, orange!80!black, domain=0:60] {80 - (-1/45*x^2 + 4/3*x)};

        % Đường gióng nét đứt từ điểm I(30; 20) sang trục tung
        \draw[dashed, thick] (axis cs:0,20) -- (axis cs:30,20);
        
        % Các điểm đặc biệt O, I, A
        \fill[red] (axis cs:0,0) circle (1.5pt) node[below left] {$O$};
        \fill[red] (axis cs:30,20) circle (1.5pt) node[above] {$I$};
        \fill[red] (axis cs:60,0) circle (1.5pt) node[above right] {$A$};
        
        % Mũi tên kích thước 20m phía dưới (từ trục x = 0 đến đỉnh parabol y = 20)
        \draw[<->, thick, gray!60!black] (axis cs:30, 0) -- (axis cs:30, 20);
        \node[fill=white, inner sep=1pt] at (axis cs:30, 10) {\small $20\text{ m}$};

        % Mũi tên kích thước 20m phía trên (từ đỉnh parabol y = 60 đến mép trên hình chữ nhật y = 80)
        \draw[<->, thick, gray!60!black] (axis cs:30, 60) -- (axis cs:30, 80);
        \node[fill=white, inner sep=1pt] at (axis cs:30, 70) {\small $20\text{ m}$};
    \end{axis}
\end{tikzpicture}
\end{center}

Nhận thấy hai phần trồng hoa dạng đường parabol có diện tích như nhau. Vậy ta xét đường parabol ở bồn hoa phía dưới có dạng $(P): y = ax^2 + bx + c$ với $a \neq 0$.

Ta có đồ thị $(P)$ đi qua ba điểm là $O(0; 0)$, $I(30; 20)$ và $A(60; 0)$. Thay vào ta thu được hệ phương trình:
\[
\begin{cases}
c = 0 \\
900a + 30b + c = 20 \\
3600a + 60b + c = 0
\end{cases}
\iff
\begin{cases}
a = -\dfrac{1}{45} \\[6pt]
b = \dfrac{4}{3} \\[6pt]
c = 0
\end{cases}
\]

Như vậy, ta có phương trình parabol đó là $y = -\dfrac{1}{45}x^2 + \dfrac{4}{3}x$ và diện tích của phần trồng hoa phía dưới là:
\[
S_{\text{hoa}} = \int_{0}^{60} \left| -\frac{1}{45}x^2 + \frac{4}{3}x \right| \mathrm{d}x = \int_{0}^{60} \left(-\frac{1}{45}x^2 + \frac{4}{3}x\right) \mathrm{d}x = 800 \, (\text{m}^2).
\]

Khi đó, diện tích của phần sân chơi là:
\[
S = 60 \cdot 80 - 2S_{\text{hoa}} = 4800 - 2 \cdot 800 = 3200 \, (\text{m}^2).
\]

\noindent\textbf{Bài 5.} Một vật trang trí có dạng một khối tròn xoay được tạo thành khi quay miền $(R)$ (phần gạch chéo trong hình vẽ) quanh trục $MN$. Biết rằng $ABCD$ là hình chữ nhật với $AB = 6\,\text{cm}, AD = 10\,\text{cm}$ và $M, N$ lần lượt là trung điểm của $AB, CD$. Hai đường cong là đường Elip có hình chữ nhật cơ sở là $ABCD$ và đường tròn tiếp xúc với hai cạnh $AD$ và $BC$ (tham khảo hình vẽ). Tính thể tích của vật trang trí đó (kết quả làm tròn đến hàng phần chục theo đơn vị $\text{cm}^3$).

\begin{center}
\begin{tikzpicture}
    \begin{axis}[
        xmin=-6.5, xmax=6.5,
        ymin=-4.2, ymax=4.2,
        axis lines=none, % Tắt hoàn toàn trục tọa độ
        ticks=none,      % Tắt các vạch chia
        samples=200,
        width=9cm, height=6.5cm,
    ]
        % Hình chữ nhật ABCD cơ sở: x từ -5 đến 5, y từ -3 đến 3
        \addplot[thick, black] coordinates {(-5,3) (5,3) (5,-3) (-5,-3) (-5,3)};
        
        % Nhãn các đỉnh hình chữ nhật A, B, C, D và M, N
        \node at (axis cs:-5, 3.3) [black] {$A$};
        \node at (axis cs:5, 3.3) [black] {$D$};
        \node at (axis cs:-5, -3.3) [black] {$B$};
        \node at (axis cs:5, -3.3) [black] {$C$};
        \node at (axis cs:-5, 0) [left, black] {$M$};
        \node at (axis cs:5, 0) [right, black] {$N$};

        % Gạch chéo toàn bộ phần bên trong hình elip
        \addplot[pattern=north east lines, pattern color=black, draw=none, domain=-5:5] {sqrt(9*(1 - x^2/25))} \closedcycle;
        \addplot[pattern=north east lines, pattern color=black, draw=none, domain=-5:5] {-sqrt(9*(1 - x^2/25))} \closedcycle;

        % Tô nền trắng bên trong hình tròn để che phần gạch chéo ở vùng giữa
        \addplot[fill=white, draw=none, domain=-3:3] {sqrt(9 - x^2)} \closedcycle;
        \addplot[fill=white, draw=none, domain=-3:3] {-sqrt(9 - x^2)} \closedcycle;

        % Vẽ viền đường Elip
        \addplot[thick, black, domain=-5:5] {sqrt(9*(1 - x^2/25))};
        \addplot[thick, black, domain=-5:5] {-sqrt(9*(1 - x^2/25))};
        
        % Vẽ viền đường tròn (C) ở giữa
        \addplot[thick, black, domain=-3:3] {sqrt(9 - x^2)};
        \addplot[thick, black, domain=-3:3] {-sqrt(9 - x^2)};
    \end{axis}
\end{tikzpicture}
\end{center}

\noindent\textbf{Lời giải:}

\noindent\textbf{Đáp án:} $75$.

\begin{center}
\begin{tikzpicture}
    \begin{axis}[
        xmin=-7, xmax=7,
        ymin=-5, ymax=5,
        axis lines=middle,
        xlabel=$x$, ylabel=$y$,
        xtick=\empty, ytick=\empty, % Tắt tick tự động để chủ động định vị bằng tay
        samples=200,
        width=10cm, height=7.5cm,
    ]
        % Hình chữ nhật ABCD cơ sở (màu đỏ)
        \addplot[thick, red] coordinates {(-5,3) (5,3) (5,-3) (-5,-3) (-5,3)};
        
        % Nhãn các đỉnh hình chữ nhật và điểm M, N, O
        \node at (axis cs:-5, 3) [above left, black] {$A$};
        \node at (axis cs:5, 3) [above right, black] {$D$};
        \node at (axis cs:-5, -3) [below left, black] {$B$};
        \node at (axis cs:5, -3) [below right, black] {$C$};
        \node at (axis cs:-5, 0) [above left, yshift=2pt] {$M$};
        \node at (axis cs:5, 0) [above right, yshift=2pt] {$N$};
        \node at (axis cs:0, 0) [above right] {$O$};

        % --- CÁC MỐC TỌA ĐỘ ĐƯỢC ĐỊNH VỊ THỦ CÔNG KHÔNG CHỒNG LÊN TRỤC VÀ ĐƯỜNG NÉT ---
        % Các số bên trái trục Oy (xích sang trái)
        \node[fill=white, inner sep=1pt, font=\small] at (axis cs:-5, 0) [below left, xshift=-4pt] {$-5$};
        \node[fill=white, inner sep=1pt, font=\small] at (axis cs:-3, 0) [below, xshift=-8pt] {$-3$};
        
        % Các số bên phải trục Oy (xích sang phải)
        \node[fill=white, inner sep=1pt, font=\small] at (axis cs:3, 0) [below, xshift=8pt] {$3$};
        \node[fill=white, inner sep=1pt, font=\small] at (axis cs:5, 0) [below right, xshift=4pt] {$5$};

        % Số 3 phía trên (xích lên cao) và số -3 phía dưới (xích xuống thấp)
        \node[fill=white, inner sep=1pt, font=\small] at (axis cs:0, 3) [above, yshift=6pt] {$3$};
        \node[fill=white, inner sep=1pt, font=\small] at (axis cs:0, -3) [below, yshift=-6pt] {$-3$};
        -------------------------------------------------------------------------------

        % Vẽ viền đường Elip
        \addplot[thick, black, domain=-5:5] {sqrt(9*(1 - x^2/25))};
        \addplot[thick, black, domain=-5:5] {-sqrt(9*(1 - x^2/25))};

        % Vẽ đường tròn tâm O, bán kính R = 3 (liền mạch, không tô màu bên trong)
        \addplot[thick, black, domain=-3:3] {sqrt(9 - x^2)};
        \addplot[thick, black, domain=-3:3] {-sqrt(9 - x^2)};
    \end{axis}
\end{tikzpicture}
\end{center}

\noindent Gắn vào hình chữ nhật $ABCD$ hệ trục tọa độ $Oxy$ với gốc $O$ trùng với giao hai đường chéo, các điểm $M, N$ nằm trên $Ox$ (tham khảo hình vẽ).

\noindent Do $MN$ là trục đối xứng của miền $(R)$ nên khối tròn xoay được tạo thành khi quay miền $(R)$ quanh trục $MN$ trùng với khối tròn xoay được tạo thành khi quay phần $(R)$ nằm phía trên $MN$ quanh $MN$.

\noindent Gọi $V_1, V_2, V$ tương ứng là thể tích của các khối tròn xoay được tạo thành khi quay Elip $(E)$, hình tròn $(C)$ và miền $(R)$ quanh $MN$ (hay quanh $Ox$), ta có $V = V_1 - V_2$.

\noindent Elip $(E)$ có các bán trục $a = \dfrac{1}{2}AD = 5,\, b = \dfrac{1}{2}AB = 3$ nên có phương trình $\dfrac{x^2}{25} + \dfrac{y^2}{9} = 1$.

\noindent Ta có $V_1 = \pi \int_{-5}^{5} 9\left(1 - \frac{x^2}{25}\right) \mathrm{d}x = \pi \left( 9x - \frac{9}{25}x^3 \right)\Bigg|_{-5}^{5} = 60\pi \, (\text{cm}^3)$.

\noindent Hình tròn $(C)$ có tâm tại $O$, bán kính $R = 3$ khi quay quanh $MN$ tạo thành khối cầu tâm $O$, bán kính bằng $3$. Do đó $V_2 = \frac{4}{3}\pi \cdot 3^3 = 36\pi \, (\text{cm}^3)$.

\noindent Thể tích cần tìm là:
\[
V = V_1 - V_2 = 60\pi - 36\pi = 24\pi \approx 75,4 \, (\text{cm}^3) \xrightarrow{\text{làm tròn}} 75.
\]


\noindent\textbf{Bài 6.} Từ hình chữ nhật $ABCD$ có chiều dài $AB = 8\text{ cm}$ và chiều rộng $BC = 4\text{ cm}$, người ta cắt bỏ miền $(R)$ được giới hạn bởi cạnh $CD$ của hình chữ nhật và hai nửa đường thẳng parabol có chung đỉnh là trung điểm cạnh $AB$, chúng lần lượt đi qua hai đầu mút $C, D$ của hình chữ nhật đó (phần tô đậm như hình vẽ). Phần còn lại cho quay quanh trục $AB$ để tạo nên một đồ vật làm trang trí, thể tích của vật trang trí đó bằng bao nhiêu $\text{cm}^3$. Viết kết quả làm tròn đến hàng đơn vị.

\begin{center}
\begin{tikzpicture}[scale=1.5, line width=0.8pt]
    % Định nghĩa các điểm
    \coordinate (A) at (0, 0);
    \coordinate (B) at (4, 0);
    \coordinate (O) at (2, 0);
    \coordinate (C) at (4, 2);
    \coordinate (D) at (0, 2);

    % 1. Tô màu vùng xanh
    % Đi từ D (0,2) sang C (4,2), vòng theo cung CO về O (tâm B(4,0)), rồi vòng theo cung OD về D (tâm A(0,0))
    \fill[cyan!30] (D) -- (C) arc (90:180:2) arc (0:90:2) -- cycle;

    % 2. Vẽ hình chữ nhật ABCD
    \draw (A) -- (B) -- (C) -- (D) -- cycle;

    % 3. Vẽ 2 cung tròn (DO tâm A, CO tâm B)
    \draw (2,0) arc (0:90:2);    % Cung OD (tâm A)
    \draw (2,0) arc (180:90:2);  % Cung OC (tâm B)

    % 4. Nhãn các đỉnh
    \node[below left] at (A) {$A$};
    \node[below right] at (B) {$B$};
    \node[below] at (O) {$O$};
    \node[above right] at (C) {$C$};
    \node[above left] at (D) {$D$};
\end{tikzpicture}
\end{center}
\noindent\textbf{Lời giải}

\noindent \textbf{Trả lời: 201}

\noindent Phương trình Parabol $OC$ có dạng: $y^2 = ax$

\noindent Mặt khác Parabol đi qua điểm $C(4; 4) \Rightarrow 4^2 = 4a \Rightarrow a = 4$

\noindent Khi đó phương trình đường cong $OC$ là $y = \sqrt{4x}$.

\noindent Thể tích của vật thể thu được là: 
$$V = 2\left(\pi \int_{0}^{4} (\sqrt{4x})^2 \,\mathrm{d}x\right) = 64\pi \approx 201\,(\text{cm}^3)$$.

\end{document}
